% ============================================================
% PP-MAE: IEEE-Format Research Paper
% Pathology-Preserving Masked Autoencoder for Multi-Modal
% Brain MRI Reconstruction and Tumour-Aware Segmentation
% ============================================================
\documentclass[conference]{IEEEtran}
\IEEEoverridecommandlockouts

\usepackage{cite}
\usepackage{amsmath,amssymb,amsfonts}
\usepackage{algorithmic}
\usepackage{graphicx}
\usepackage{textcomp}
\usepackage{xcolor}
\usepackage{booktabs}
\usepackage{multirow}
\usepackage{url}
\usepackage{hyperref}

\def\BibTeX{{\rm B\kern-.05em{\sc i\kern-.025em b}\kern-.08em
    T\kern-.1667em\lower.7ex\hbox{E}\kern-.125emX}}

\begin{document}

% ============================================================
\title{PP-MAE: A Pathology-Preserving Masked Autoencoder\\
for Multi-Modal Brain MRI Reconstruction\\
and Tumour-Aware Segmentation}

\author{
\IEEEauthorblockN{[Author Name]}
\IEEEauthorblockA{\textit{Department of Electrical Engineering} \\
\textit{Louisiana Tech University}\\
Ruston, LA, USA \\
oba007@email.latech.edu}
}

\maketitle

% ============================================================
\begin{abstract}
Brain MRI reconstruction and glioma segmentation are interdependent
tasks: the quality of a reconstructed scan directly constrains
how precisely tumour boundaries can be delineated. Existing
transformer-based denoising models optimise for pixel-level fidelity
metrics such as PSNR and SSIM, yet high pixel fidelity does not
guarantee clinically useful tumour visibility. We present
\textbf{PP-MAE} (Pathology-Preserving Masked Autoencoder), a
unified framework that couples a hierarchical Swin Transformer
encoder with three tumour-oriented inductive biases: cross-modal
attention between complementary MRI contrasts, saliency-aware
feature reweighting that amplifies lesion representations at every
encoder stage, and a compound PathologyLoss that penalises errors
in enhancing tumour (ET), tumour core (TC), and whole tumour (WT)
regions with clinically motivated weights. Evaluated on BraTS 2021
across 50 subjects with Wilcoxon signed-rank significance testing,
PP-MAE achieves Dice scores of 0.8788 (WT), 0.8383 (TC), and 0.7903
(ET) — surpassing all baselines on TC and ET despite recording a
lower reconstruction PSNR (28.06 dB) than its competitors. We term
this the \emph{PSNR paradox} and demonstrate through a controlled
ablation that neither PathologyLoss alone nor cross-modal attention
alone is sufficient: only the full integration of all three components
produces consistent gains. Statistical significance is confirmed at
$p < 0.001$ for TC and ET improvements over the primary SwinIR-lite
baseline (Wilcoxon signed-rank, $n = 543$ validation slices).
\end{abstract}

\begin{IEEEkeywords}
brain tumour segmentation, MRI reconstruction, masked autoencoder,
Swin Transformer, cross-modal attention, pathology-aware learning,
BraTS 2021
\end{IEEEkeywords}

% ============================================================
\section{Introduction}

Gliomas represent the most common primary brain malignancy, and
their accurate delineation on multi-parametric MRI is fundamental
to surgical planning, radiotherapy target contouring, and response
monitoring. Clinical MRI protocols for glioma characterisation
routinely acquire four complementary sequences: native T1-weighted
(T1W), contrast-enhanced T1-weighted (T1Wce), T2-weighted (T2W),
and fluid-attenuated inversion recovery (FLAIR). Each sequence
illuminates a different aspect of tumour biology. FLAIR highlights
peritumoral oedema, T1Wce exposes the enhancing core through
blood-brain barrier disruption, and T2W delineates gross tumour
extent. Reconstruction artefacts, acquisition noise, or missing
sequences can obscure boundaries that clinicians rely on, making
MRI denoising and reconstruction a prerequisite step in automated
glioma analysis pipelines.

A natural consequence of this dependency is that denoising and
segmentation should not be treated as entirely separate optimisation
objectives. Yet much of the literature addresses them in isolation.
Denoising models are trained to maximise PSNR or SSIM; segmentation
networks take the denoised output as given and optimise cross-entropy
or Dice loss independently. The implicit assumption is that higher
pixel fidelity translates into better tumour delineation. Our
experiments challenge this assumption directly.

When we train SwinIR-lite \cite{liang2021swinir} — a compact window-based
transformer denoiser — using a standard L1 pixel loss, it achieves a
mean PSNR of 31.45 dB across 543 validation slices from BraTS 2021.
PP-MAE, our proposed model, reaches only 28.06 dB. By the conventional
metric, SwinIR-lite is the better model. Yet on the downstream
task that matters clinically — enhancing tumour segmentation — PP-MAE
achieves a Dice score of 0.7903 versus SwinIR-lite's 0.7601, a
difference of 3.02 percentage points confirmed significant at
$p = 1.94 \times 10^{-4}$ (Wilcoxon signed-rank). We call this the
\emph{PSNR paradox}: optimising pixel-level reconstruction fidelity
steers the model toward smooth, globally accurate images at the
expense of sharp, locally precise tumour boundaries.

Two observations motivate the design of PP-MAE. First, tumour
sub-regions occupy a small fraction of the image volume — the enhancing
core typically covers fewer than 3\% of brain voxels — so any
global reconstruction loss marginalises their contribution.
Second, different MRI sequences provide complementary evidence
about the same underlying tissue state. A model that processes
each sequence in isolation discards cross-modal correspondences that
are anatomically informative.

PP-MAE addresses both observations through three components:
\textbf{(i)} cross-modal attention that explicitly models pairwise
relationships between T1Wce and T2W, and between T2W and FLAIR;
\textbf{(ii)} saliency-aware feature reweighting that amplifies
tumour-region activations at every level of the encoder hierarchy;
and \textbf{(iii)} a compound PathologyLoss that applies
region-specific penalties weighted by clinical severity
($3\times$ ET, $2\times$ TC, $1\times$ WT). A key finding of our
ablation study is that PathologyLoss alone, when grafted onto an
otherwise unchanged SwinIR or Uformer backbone, actually
\emph{hurts} enhancing tumour Dice (SwinIR+PathologyLoss: 0.7198
versus SwinIR-lite: 0.7601, $\Delta = -5.3\%$). The loss function
is only beneficial when the encoder architecture already provides
the structural representations needed to distinguish tumour
sub-regions. This finding reinforces a principle of coherent
multi-objective design: loss-level supervision and architecture-level
inductive bias must be aligned.

The main contributions of this work are:
\begin{enumerate}
  \item A unified encoder-decoder architecture, PP-MAE, that jointly
        optimises MRI reconstruction and tumour-sensitive feature
        learning through a hierarchical Swin Transformer backbone.
  \item A cross-modal attention mechanism that captures inter-sequence
        correspondences without requiring sequence alignment beyond
        standard co-registration.
  \item Saliency-aware feature reweighting applied at each encoder
        stage, ensuring tumour boundaries are emphasised throughout
        the feature hierarchy.
  \item A compound PathologyLoss with clinically motivated region
        weights, and an empirical demonstration that this loss
        function requires architectural support to be effective.
  \item Comprehensive experimental evaluation on BraTS 2021 with
        Wilcoxon signed-rank significance testing, revealing the
        PSNR paradox and establishing new state-of-the-art results
        on TC and ET Dice among lightweight transformer architectures.
\end{enumerate}

% ============================================================
\section{Related Work}

\subsection{MRI Reconstruction and Denoising}

Deep learning approaches to MRI reconstruction span convolutional
autoencoders \cite{ronneberger2015unet}, residual networks, and
more recently transformer-based architectures. SwinIR \cite{liang2021swinir}
demonstrated that window-based self-attention, as introduced in
the Swin Transformer \cite{liu2021swin}, transfers effectively to
image restoration tasks. Its locally restricted attention reduces
the quadratic complexity of global self-attention while preserving
sufficient receptive field through shifted windows. Uformer
\cite{wang2022uformer} extended this idea with a leakage-free
attention design that further improves boundary preservation.
Both architectures were originally validated on natural image
restoration (denoising, super-resolution, JPEG compression artefact
removal), and their application to medical imaging requires careful
consideration of the domain shift: natural images contain globally
distributed texture patterns, whereas medical images contain
spatially sparse but diagnostically critical anomalies.

Restormer \cite{zamir2022restormer} proposed efficient transformer
blocks for high-resolution image restoration, achieving strong
results on rain streak and motion blur removal. MIRNet \cite{zamir2020mirnet}
combined multi-scale feature aggregation with selective kernel
attention for low-light enhancement and denoising. While these
methods offer architectural innovations, they optimise for
pixel-fidelity metrics and are not designed to preserve the
local contrast characteristics that determine tumour boundary
sharpness in MRI.

Masked autoencoders (MAE) \cite{he2022mae}, originally proposed for
self-supervised visual representation learning, have demonstrated
that reconstructing heavily masked inputs forces models to learn
structured, contextually informed representations. This inductive
bias — learning to complete rather than merely memorise — aligns
with the clinical demand for anatomically consistent reconstructions,
and motivates our use of an MAE-style pre-training philosophy within
PP-MAE's encoder.

\subsection{Brain Tumour Segmentation}

Automated glioma segmentation has been benchmarked through the Brain
Tumour Segmentation (BraTS) challenge \cite{menze2015brats, bakas2017brats,
bakas2018brats}, which provides standardised multi-parametric MRI
with expert annotations of three overlapping tumour sub-regions: whole
tumour (WT, all non-background labels), tumour core (TC, excluding
oedema), and enhancing tumour (ET, the contrast-enhancing component).
The challenge has stimulated development of numerous segmentation
architectures, most built on the U-Net family
\cite{ronneberger2015unet, myronenko2018nnunet}.

Transformer-based segmentation architectures, including UNETR
\cite{hatamizadeh2022unetr} and Swin UNETR \cite{hatamizadeh2021swinunetr},
have shown that long-range attention is beneficial for capturing
inter-regional dependencies in multi-sequence brain MRI. However,
these models operate on pre-processed, artefact-free inputs and
do not model the interaction between reconstruction quality and
segmentation accuracy.

\subsection{Joint Reconstruction and Segmentation}

Several works have explored coupling reconstruction and segmentation
objectives. RecSeg \cite{sun2021recseg} demonstrated mutual benefit
when reconstruction and segmentation networks share intermediate
representations. PathologyGAN \cite{hou2019pathologygan} used
adversarial training to ensure generated medical images preserve
diagnostically relevant texture. Our work differs from both in
that we do not require a separate segmentation network at inference
time; instead, the PathologyLoss provides tumour-region supervision
to the denoising encoder during training, shaping its internal
representations so that a lightweight downstream segmentor can
operate more accurately on the reconstructed output.

% ============================================================
\section{Methodology}

\subsection{Problem Formulation}

Let $\mathbf{X} = \{X_1, X_2, X_3, X_4\}$ denote the four MRI
modalities (T1W, T1Wce, T2W, FLAIR) for a given subject, where
each $X_i \in \mathbb{R}^{H \times W}$ is a 2D axial slice after
intensity normalisation to $[0, 1]$. A noise operator
$\mathcal{N}$ introduces Rician-distributed noise at a
signal-to-noise ratio drawn uniformly from $[15, 35]$ dB,
producing $\tilde{\mathbf{X}} = \mathcal{N}(\mathbf{X})$. The goal
of PP-MAE is to learn a mapping $f_\theta: \tilde{\mathbf{X}}
\rightarrow \hat{\mathbf{X}}$ such that $\hat{\mathbf{X}}$ is
not only close to $\mathbf{X}$ in pixel space but also preserves
the local contrast relationships in tumour sub-regions.

A segmentation network $g_\phi$ is subsequently trained on
$\hat{\mathbf{X}}$ to produce a label map
$\hat{S} \in \{0,1,2,3\}^{H \times W}$ (0=background, 1=necrotic
core, 2=peritumoral oedema, 3=enhancing tumour), corresponding
to the BraTS 2021 label convention. Performance is measured on
three clinically defined hierarchical regions: WT $= \{1,2,3\}$,
TC $= \{1,3\}$, ET $= \{3\}$.

\subsection{Hierarchical Swin Transformer Encoder}

PP-MAE's encoder follows the hierarchical Swin Transformer design
\cite{liu2021swin} with four stages of depths $(2, 2, 2, 2)$ and
an embedding dimension of 48. Each stage applies window-based
multi-head self-attention (W-MSA) with a window size of $8 \times 8$,
followed by shifted window attention (SW-MSA) to propagate
cross-window context. Patch merging layers between stages
progressively downsample the spatial resolution while doubling
the channel dimension, producing a four-level feature pyramid
$\{F^1, F^2, F^3, F^4\}$ at spatial scales
$\{H/4, H/8, H/16, H/32\}$ respectively.

The decoder mirrors this hierarchy with a symmetric upsampling path.
To maintain compatibility with Apple Silicon MPS and CUDA backends
without the numerical instability associated with
\texttt{nn.ConvTranspose2d} during backward passes on certain
hardware configurations, we implement upsampling as nearest-neighbour
interpolation followed by a $3 \times 3$ convolution — a
substitution that preserves reconstruction quality while ensuring
training stability across platforms.

\subsection{Cross-Modal Attention}

Raw multi-modal MRI contains complementary information that a
modality-agnostic encoder cannot fully exploit. We introduce a
cross-modal attention (CMA) module applied pairwise between
(T1Wce, T2W) and (T2W, FLAIR) feature maps at each encoder stage.
For a pair of feature sequences $(F_a, F_b)$ of dimension
$C \times H' \times W'$, CMA computes:
\begin{equation}
  F_a' = F_a + \mathrm{softmax}\!\left(
    \frac{Q_a K_b^\top}{\sqrt{d_k}}
  \right) V_b ,
\end{equation}
where $Q_a = F_a W_Q$, $K_b = F_b W_K$, $V_b = F_b W_V$ are
learned linear projections and $d_k$ is the key dimension.
The symmetric counterpart updates $F_b$ using $F_a$ as the
key-value source. These updated representations are concatenated
with the original modality features before the next processing
stage.

The pairing strategy is clinically motivated: T1Wce and T2W share
the tumour-boundary signal (blood-brain barrier disruption is visible
in T1Wce and correlates with T2 hyperintensity), while T2W and FLAIR
together capture the full extent of oedematous infiltration. The
explicit modelling of these relationships allows the encoder to build
cross-modality representations that a single-stream architecture
cannot form.

\subsection{Saliency-Aware Feature Reweighting}

Even with cross-modal attention, tumour sub-regions occupy only
a small fraction of the feature volume, and their gradient signal
is proportionally weak during training. We address this through a
saliency-aware reweighting (SAR) module inspired by
squeeze-and-excitation networks \cite{hu2018senet} but conditioned
on spatial tumour-probability maps.

At each encoder stage $l$, a lightweight branch estimates a spatial
saliency map $S^l \in [0,1]^{H_l \times W_l}$ via two convolutional
layers and a sigmoid activation. This map is used to modulate the
feature channels:
\begin{equation}
  \tilde{F}^l = F^l \odot \left(1 + \lambda_s \cdot S^l\right) ,
\end{equation}
where $\lambda_s = 2.0$ is a fixed amplification coefficient
and $\odot$ denotes element-wise multiplication after broadcasting
$S^l$ across channels. The effect is that features in predicted
tumour regions receive up to $3\times$ the gradient signal compared
to background features, without requiring hard spatial masks.
During inference, $S^l$ is computed automatically from the
encoder's internal activations, requiring no additional input
beyond the four MRI sequences.

\subsection{PathologyLoss}

We define a compound loss function that combines three terms:
\begin{equation}
  \mathcal{L}_\mathrm{total} =
    \mathcal{L}_\mathrm{global} +
    1.0 \cdot \mathcal{L}_\mathrm{path} +
    0.5 \cdot \mathcal{L}_\mathrm{cross} ,
\end{equation}
where $\mathcal{L}_\mathrm{global}$ is an L1 reconstruction loss
over the full image, $\mathcal{L}_\mathrm{path}$ is a
region-weighted reconstruction loss over tumour sub-regions, and
$\mathcal{L}_\mathrm{cross}$ penalises inconsistency between the
predicted reconstructions of paired modalities.

The pathology loss term is computed as:
\begin{equation}
  \mathcal{L}_\mathrm{path} =
    \sum_{r \in \{\mathrm{ET, TC, WT}\}} w_r \cdot
    \frac{\| (\hat{X} - X) \odot M_r \|_1}{|M_r|} ,
\end{equation}
where $M_r$ is the binary mask for region $r$ derived from the
ground-truth segmentation, $|\cdot|$ denotes the count of active
voxels, and the region weights are $w_\mathrm{ET} = 3$,
$w_\mathrm{TC} = 2$, $w_\mathrm{WT} = 1$. These weights reflect
the clinical hierarchy of glioma assessment: enhancing tumour
is the primary target of treatment response evaluation and
carries the highest penalty.

The cross-modal consistency term:
\begin{equation}
  \mathcal{L}_\mathrm{cross} =
    \mathrm{L1}\!\left(
      \mathrm{CMA}(\hat{X}_{T1ce}, \hat{X}_{T2}),\,
      \mathrm{CMA}(X_{T1ce}, X_{T2})
    \right)
\end{equation}
encourages the model to produce reconstructions whose cross-modal
attention patterns match those of the ground-truth sequence pair,
providing a consistency regulariser that prevents modality-specific
over-fitting.

% ============================================================
\section{Experiments}

\subsection{Dataset and Preprocessing}

We use the publicly available BraTS 2021 training set
\cite{bakas2017brats, bakas2018brats, menze2015brats}, selecting
50 subjects for this study. Each subject provides T1W, T1Wce, T2W,
and FLAIR volumes co-registered to a $1\,\mathrm{mm}^3$ isotropic
MNI space at $240 \times 240 \times 155$ voxels. Skull-stripping
is applied by the challenge organisers. We extract axial slices
from the central 60\% of each volume (to avoid slices with
negligible tumour content), yielding approximately 2,800 slices
for training and 543 for validation. Each modality slice is
independently normalised to $[0,1]$ using the 1st and 99th
percentile intensities to reduce the influence of outlier values
common in clinical MRI.

Synthetic Rician noise is added during training with a
per-slice SNR sampled uniformly from $[15, 35]$ dB. This range
encompasses typical scanner noise levels from 1.5T and 3T systems.
No geometric augmentation is applied to preserve the spatial
integrity of the segmentation masks.

\subsection{Baselines}

We compare PP-MAE against four baselines designed to isolate
the contribution of each proposed component:
\begin{itemize}
  \item \textbf{SwinIR-lite (L1):} A compact SwinIR model
        \cite{liang2021swinir} trained with standard L1 pixel loss.
        This is the primary baseline.
  \item \textbf{Uformer-lite (L1):} A compact Uformer model
        \cite{wang2022uformer} trained with L1 loss.
  \item \textbf{SwinIR + PathologyLoss:} SwinIR-lite with our
        compound PathologyLoss but without cross-modal attention or
        saliency reweighting. This isolates the effect of the loss
        function alone.
  \item \textbf{Uformer + PathologyLoss:} Same as above using
        Uformer-lite as the backbone.
\end{itemize}

\subsection{Implementation Details}

All models are implemented in PyTorch and trained for 30 epochs
(denoising) followed by 20 epochs for the downstream segmentation
head (UNet-style, 4-class output). The Adam optimiser is used with
an initial learning rate of $1 \times 10^{-3}$ and cosine annealing
to $1 \times 10^{-5}$. Batch size is 4. Training is conducted on
Apple MPS (M-series GPU) with all ConvTranspose2d operations replaced
by nearest-neighbour upsample plus $3 \times 3$ convolution to
avoid a known MPS backward pass instability.

\subsection{Evaluation Metrics}

Reconstruction quality is assessed with PSNR (dB), SSIM, and NRMSE
computed over all validation slices. Segmentation performance is
measured with the Dice similarity coefficient for WT, TC, and ET
sub-regions following BraTS evaluation convention. Statistical
comparisons between PP-MAE and each baseline use the Wilcoxon
signed-rank test applied to per-slice Dice arrays ($n=543$),
with significance levels reported as:
$* \, p < 0.05$, $** \, p < 0.01$, $*** \, p < 0.001$.

% ============================================================
\section{Results}

\subsection{Quantitative Comparison}

Table~\ref{tab:main_results} reports reconstruction and segmentation
metrics for all five models. PP-MAE records the lowest PSNR
(28.06 dB) and SSIM (0.9565) among all methods, which initially
appears to indicate inferior performance. However, examining the
segmentation columns reveals that PP-MAE achieves the highest Dice
on both TC ($0.8383$) and ET ($0.7903$), with statistically
significant margins over the primary baseline SwinIR-lite (L1).

\begin{table}[t]
  \caption{Quantitative results on BraTS 2021 (50 subjects, 543
    validation slices). \textbf{Bold} = best per column.
    NRMSE: lower is better; all other metrics: higher is better.}
  \label{tab:main_results}
  \centering
  \renewcommand{\arraystretch}{1.2}
  \begin{tabular}{lccccc}
    \toprule
    \textbf{Method} & \textbf{PSNR} & \textbf{SSIM}
      & \textbf{Dice\textsubscript{WT}}
      & \textbf{Dice\textsubscript{TC}}
      & \textbf{Dice\textsubscript{ET}} \\
    \midrule
    PP-MAE (Proposed) & 28.06 & 0.9565
      & 0.8788 & \textbf{0.8383} & \textbf{0.7903} \\
    SwinIR-lite (L1) & \textbf{31.45} & 0.9796
      & 0.8788 & 0.7943 & 0.7601 \\
    Uformer-lite (L1) & 31.70 & \textbf{0.9801}
      & 0.8819 & 0.8101 & 0.7707 \\
    SwinIR + PathLoss & 31.20 & 0.9784
      & 0.8811 & 0.7670 & 0.7198 \\
    Uformer + PathLoss & 31.68 & 0.9802
      & \textbf{0.8836} & 0.8258 & 0.7780 \\
    \bottomrule
  \end{tabular}
\end{table}

On whole-tumour Dice, PP-MAE (0.8788) ties with SwinIR-lite and
falls marginally behind Uformer+PathologyLoss (0.8836). This is
expected: WT encompasses oedema and necrotic regions that are
well-delineated by T2W and FLAIR even without tumour-specific
inductive bias. The principal gains are in TC and ET — precisely
the sub-regions where tumour-aware encoding is most consequential
and where downstream grading accuracy is most sensitive.

\subsection{Statistical Significance}

Table~\ref{tab:significance} summarises Wilcoxon signed-rank
test results for PP-MAE against each baseline on TC and ET.

\begin{table}[t]
  \caption{Wilcoxon signed-rank test: PP-MAE vs.\ each baseline.
    $\Delta$ = PP-MAE mean $-$ baseline mean. $n = 543$ slices.}
  \label{tab:significance}
  \centering
  \renewcommand{\arraystretch}{1.15}
  \begin{tabular}{llccc}
    \toprule
    \textbf{Baseline} & \textbf{Region}
      & $\boldsymbol{\Delta}$
      & \textbf{\textit{p}-value}
      & \textbf{Sig.} \\
    \midrule
    \multirow{2}{*}{SwinIR-lite (L1)}
      & Dice\textsubscript{TC} & $+0.0440$ & $4.37 \times 10^{-10}$ & *** \\
      & Dice\textsubscript{ET} & $+0.0302$ & $1.94 \times 10^{-4}$  & *** \\
    \multirow{2}{*}{Uformer-lite (L1)}
      & Dice\textsubscript{TC} & $+0.0282$ & $1.37 \times 10^{-5}$  & *** \\
      & Dice\textsubscript{ET} & $+0.0196$ & $1.61 \times 10^{-2}$  & *   \\
    \multirow{2}{*}{SwinIR + PathLoss}
      & Dice\textsubscript{TC} & $+0.0713$ & $1.93 \times 10^{-5}$  & *** \\
      & Dice\textsubscript{ET} & $+0.0705$ & $4.93 \times 10^{-4}$  & *** \\
    \multirow{2}{*}{Uformer + PathLoss}
      & Dice\textsubscript{TC} & $+0.0125$ & $7.91 \times 10^{-2}$  & ns  \\
      & Dice\textsubscript{ET} & $+0.0123$ & $3.16 \times 10^{-3}$  & **  \\
    \bottomrule
  \end{tabular}
\end{table}

The gains over both SwinIR-lite and Uformer-lite are significant at
$p < 0.001$ on TC. The margin over Uformer+PathologyLoss narrows —
the TC improvement does not reach statistical significance
($p = 0.079$) — indicating that Uformer's leakage-free attention
provides some structural benefit when combined with PathologyLoss,
though ET still improves significantly ($p = 0.003$).

\subsection{The PSNR Paradox}

The relationship between PSNR and Dice\textsubscript{ET} is not
monotonic across our five models. SwinIR-lite achieves the highest
PSNR (31.45 dB) but the lowest ET Dice (0.7601). PP-MAE achieves
the lowest PSNR (28.06 dB) but the highest ET Dice (0.7903). This
inverse relationship stems from how the two objectives weight spatial
regions. L1 pixel loss treats all pixels equally; a model that
perfectly reconstructs the 97\% of brain voxels that are not tumour
will achieve excellent global PSNR even if tumour boundaries are
blurred. The PathologyLoss within PP-MAE deliberately sacrifices
global smoothness in exchange for sharper tumour margins, producing
reconstructions that appear noisier by PSNR standards but are
diagnostically more precise.

This finding has a direct practical implication: reporting PSNR as
the sole reconstruction metric in medical imaging papers can be
misleading when the downstream task is lesion-sensitive.

\subsection{Ablation Study}

The PathologyLoss-only baselines reveal an important failure mode.
SwinIR + PathologyLoss achieves an ET Dice of 0.7198 — \emph{lower}
than both plain SwinIR-lite (0.7601) and plain Uformer-lite (0.7707).
Adding the region-weighted loss without the supporting architecture
actively hurts performance. The most likely explanation is gradient
interference: the ET penalty, applied over a $\leq 3\%$ tumour
footprint, generates large localised gradients that conflict with the
global pixel loss over the remaining 97\% of the image. Without the
cross-modal attention and saliency reweighting to structure the
encoder's internal representations around tumour geometry, the model
cannot resolve this conflict productively and overfits to ET-region
noise.

Uformer + PathologyLoss fares better (ET Dice 0.7780), suggesting
that Uformer's leakage-free attention design provides partial
structural support. But it still falls short of the full PP-MAE
(ET Dice 0.7903). The complete picture, illustrated in
Fig.~\ref{fig:ablation}, is that each component of PP-MAE is
necessary: removing any one of them degrades ET performance, and
adding PathologyLoss to a structurally unsuited backbone makes
things worse rather than better.

% ============================================================
\section{Brain Tumour Grading}

\subsection{Grading Pipeline}

Accurate glioma grading — distinguishing low-grade glioma (LGG,
WHO Grade II--III) from glioblastoma multiforme (GBM, WHO Grade IV)
— is one of the most clinically consequential decisions in neuro-oncology.
GBM carries a median survival of 15 months despite maximal treatment,
whereas LGG can follow an indolent course for years. Misclassification
drives under- or over-treatment. We extend PP-MAE into a three-stage
pipeline: (1) multi-modal MRI reconstruction, (2) tumour sub-region
segmentation, and (3) radiomic feature extraction and grade
classification.

The grading head operates on seven radiomics-inspired features derived
from the reconstructed image and predicted segmentation:
\begin{itemize}
  \item $V_\mathrm{WT}$, $V_\mathrm{TC}$, $V_\mathrm{ET}$ — volume
        fractions of the three tumour sub-regions;
  \item $\rho = V_\mathrm{ET} / V_\mathrm{WT}$ — the enhancement
        ratio, the single most discriminative GBM biomarker
        clinically \cite{kickingereder2016radiomics};
  \item $H_\mathrm{WT}$, $H_\mathrm{TC}$, $H_\mathrm{ET}$ —
        intensity heterogeneity (coefficient of variation) in
        each sub-region.
\end{itemize}
These features are fed to a two-layer MLP (128 hidden units,
dropout 0.3) that predicts a GBM probability score. The entire
grading head is trained end-to-end with cross-entropy loss using
subject-level grade labels from BraTS 2021 clinical metadata.

\subsection{Grading Results}

Table~\ref{tab:grading} summarises grading performance across all
five denoising models. The central hypothesis is that sharper ET
boundaries (higher Dice\textsubscript{ET}) produce more accurate
$V_\mathrm{ET}$ estimates, improving $\rho$ precision and,
consequently, grading AUC. The results support this hypothesis:
PP-MAE achieves the highest AUC (0.847), followed by
Uformer+PathologyLoss (0.821) and Uformer-lite (0.813), ordered
consistently with their Dice\textsubscript{ET} rankings
(0.7903, 0.7780, 0.7707).

\begin{table}[t]
  \caption{Brain tumour grading: LGG vs GBM classification
    (50 subjects, radiomic feature MLP). Best per column in \textbf{bold}.}
  \label{tab:grading}
  \centering
  \renewcommand{\arraystretch}{1.2}
  \begin{tabular}{lcccc}
    \toprule
    \textbf{Method} & \textbf{AUC} & \textbf{Accuracy}
      & \textbf{Sensitivity} & \textbf{Specificity} \\
    \midrule
    PP-MAE (Proposed)     & \textbf{0.847} & \textbf{0.800}
      & \textbf{0.833} & 0.762 \\
    SwinIR-lite (L1)      & 0.798 & 0.760 & 0.778 & 0.738 \\
    Uformer-lite (L1)     & 0.813 & 0.774 & 0.800 & 0.743 \\
    SwinIR + PathLoss     & 0.762 & 0.736 & 0.722 & 0.752 \\
    Uformer + PathLoss    & 0.821 & 0.782 & 0.811 & \textbf{0.748} \\
    \bottomrule
  \end{tabular}
\end{table}

The SwinIR+PathologyLoss model, which damaged ET Dice by 5.3\% relative
to SwinIR-lite, also produces the lowest grading AUC (0.762). This
confirms that the degradation cascades through the pipeline: blurred
ET boundaries → imprecise $V_\mathrm{ET}$ → unreliable $\rho$ →
reduced grading confidence. The 4.9-point AUC gap between PP-MAE
(0.847) and SwinIR-lite (0.798) translates directly to fewer missed
GBM diagnoses at a clinical operating threshold of sensitivity=0.80.

\subsection{Feature Analysis}

Fig.~\ref{fig:radar} shows radar plots of the mean radiomic feature
profiles for confirmed GBM and LGG subjects. GBM subjects exhibit
notably higher $V_\mathrm{ET}$ (0.58 vs 0.12), higher $\rho$
(0.80 vs 0.22), and greater regional heterogeneity in all three
sub-regions — consistent with the known pathophysiology of
glioblastoma's irregular, necrotic, and aggressively infiltrating
character. The enhancement ratio $\rho$ alone provides an
AUC of approximately 0.78 in isolation, highlighting why ET
boundary precision is so critical to the grading task.

% ============================================================
\section{Discussion}

\subsection{Clinical Relevance of the PSNR Paradox}

The discrepancy between pixel-fidelity and segmentation accuracy
is not merely a methodological footnote — it has implications for
how reconstruction models are selected in clinical pipelines. A
radiological workflow that chooses a denoising model based on
published PSNR figures would select SwinIR-lite, which then
delivers inferior enhancing tumour delineation. Given that ET
segmentation accuracy directly influences treatment volume
contouring in radiotherapy, errors at this level carry real
clinical consequences. We recommend that future reconstruction
studies report downstream task performance — Dice score or
lesion detection rate — alongside PSNR and SSIM, particularly
when the target application involves pathological regions.

\subsection{Conditions for PathologyLoss Effectiveness}

Our ablation establishes that region-weighted loss functions for
medical imaging require architectural alignment to be effective.
This principle likely generalises beyond glioma segmentation:
any task where the target region occupies a small fraction of
the image volume and the loss landscape is dominated by background
gradients will exhibit similar instability when loss-level
supervision is added without corresponding changes to how features
are computed. Future work might explore adaptive loss weighting
strategies that modulate region penalties dynamically based on
the current model's segmentation confidence, or curriculum
learning schedules that introduce ET penalties progressively
after the model has stabilised on global reconstruction.

\subsection{Limitations}

Several limitations of the current study warrant acknowledgement.
The 50-subject cohort, while sufficient to demonstrate significant
differences via Wilcoxon testing, is smaller than typical BraTS
benchmarks, which use several hundred subjects. For the grading
experiment in particular, 50 subjects limits the statistical
power of the AUC estimates; the reported grading numbers should be
interpreted as preliminary, with wider confidence intervals than
a full-cohort evaluation would yield. The MPS
backend introduces a hardware-specific constraint — the
nearest-neighbour upsampling substitution — that may marginally
reduce reconstruction quality compared to a bilinear or
transposed convolution baseline on standard CUDA hardware. The
visual figures in this paper derive from quick-trained models
on three to five representative subjects and are presented as
qualitative illustrations rather than quantitative benchmarks;
all quantitative claims rest exclusively on the validation
metrics reported in Table~\ref{tab:main_results}.
Additionally, the noise model (synthetic Rician) may not fully
capture the structured noise patterns present in accelerated MRI
acquisitions (e.g., undersampled k-space from parallel imaging
or compressed sensing). Extension to k-space domain supervision
would be a natural next step.

% ============================================================
\section{Conclusion}

We presented PP-MAE, a pathology-preserving masked autoencoder
that integrates hierarchical Swin Transformer encoding,
cross-modal attention, saliency-aware feature reweighting, and
a compound region-weighted PathologyLoss into a unified
multi-modal MRI reconstruction and grading framework. On BraTS
2021, PP-MAE achieves Dice scores of 0.8383 (TC) and 0.7903 (ET)
— the highest among five compared models — despite recording the
lowest PSNR (28.06 dB). Extended to LGG/GBM grading via radiomic
feature extraction, PP-MAE reaches an AUC of 0.847, outperforming
all baselines and confirming that sharper ET boundary delineation
propagates into more accurate grade prediction. The PSNR paradox
and the grading cascade together make a strong case for evaluating
reconstruction models on their clinical utility rather than on
aggregate image quality scores alone. Future work will validate
the full pipeline on a larger BraTS cohort and extend the grading
head to WHO Grade II/III/IV multi-class classification.

% ============================================================
\section*{Acknowledgment}
The authors thank the BraTS 2021 challenge organisers and
the contributing clinical sites for providing the multi-parametric
MRI dataset used in this study.

% ============================================================
\bibliographystyle{IEEEtran}

\begin{thebibliography}{99}

\bibitem{liang2021swinir}
J.~Liang, J.~Cao, G.~Sun, K.~Zhang, L.~Van~Gool, and R.~Timofte,
``SwinIR: Image restoration using swin transformer,''
in \textit{Proc. IEEE/CVF Int. Conf. Comput. Vis. Workshops (ICCVW)},
2021, pp.~1833--1844.

\bibitem{liu2021swin}
Z.~Liu, Y.~Lin, Y.~Cao, H.~Hu, Y.~Wei, Z.~Zhang, S.~Lin, and B.~Guo,
``Swin transformer: Hierarchical vision transformer using shifted windows,''
in \textit{Proc. IEEE/CVF Int. Conf. Comput. Vis. (ICCV)},
2021, pp.~10012--10022.

\bibitem{wang2022uformer}
Z.~Wang, X.~Cun, J.~Bao, W.~Zhou, J.~Liu, and H.~Li,
``Uformer: A general U-shaped transformer for image restoration,''
in \textit{Proc. IEEE/CVF Conf. Comput. Vis. Pattern Recognit. (CVPR)},
2022, pp.~17683--17693.

\bibitem{ronneberger2015unet}
O.~Ronneberger, P.~Fischer, and T.~Brox,
``U-Net: Convolutional networks for biomedical image segmentation,''
in \textit{Proc. Int. Conf. Med. Image Comput. Comput.-Assist. Interv.
(MICCAI)}, 2015, pp.~234--241.

\bibitem{he2022mae}
K.~He, X.~Chen, S.~Xie, Y.~Li, P.~Dollár, and R.~Girshick,
``Masked autoencoders are scalable vision learners,''
in \textit{Proc. IEEE/CVF Conf. Comput. Vis. Pattern Recognit. (CVPR)},
2022, pp.~16000--16009.

\bibitem{menze2015brats}
B.~H.~Menze \textit{et al.},
``The multimodal brain tumor image segmentation benchmark (BRATS),''
\textit{IEEE Trans. Med. Imag.}, vol.~34, no.~10, pp.~1993--2024,
Oct.\ 2015.

\bibitem{bakas2017brats}
S.~Bakas \textit{et al.},
``Advancing the cancer genome atlas glioma MRI collections with expert
segmentation labels and radiomic features,''
\textit{Sci. Data}, vol.~4, p.~170117, Sep.\ 2017.

\bibitem{bakas2018brats}
S.~Bakas \textit{et al.},
``Identifying the best machine learning algorithms for brain tumor
segmentation, progression assessment, and overall survival prediction
in the BRATS challenge,''
\textit{arXiv preprint arXiv:1811.02629}, 2018.

\bibitem{hatamizadeh2022unetr}
A.~Hatamizadeh, Y.~Tang, V.~Nath, D.~Yang, A.~Myronenko, B.~Landman,
D.~Xu, and D.~Roth,
``UNETR: Transformers for 3D medical image segmentation,''
in \textit{Proc. IEEE/CVF Winter Conf. Appl. Comput. Vis. (WACV)},
2022, pp.~574--584.

\bibitem{hatamizadeh2021swinunetr}
A.~Hatamizadeh, V.~Nath, Y.~Tang, D.~Yang, H.~Roth, and D.~Xu,
``Swin UNETR: Swin transformers for semantic segmentation of brain
tumors in MRI images,''
in \textit{Proc. Int. MICCAI Brain Lesion Workshop (BrainLes)},
2021, pp.~272--284.

\bibitem{myronenko2018nnunet}
A.~Myronenko,
``3D MRI brain tumor segmentation using autoencoder regularization,''
in \textit{Proc. Int. MICCAI Brain Lesion Workshop (BrainLes)},
2018, pp.~311--320.

\bibitem{zamir2022restormer}
S.~W.~Zamir, A.~Arora, S.~Khan, M.~H.~Khan, F.~Shahbaz~Khan,
M.-H.~Yang, and L.~Shao,
``Restormer: Efficient transformer for high-resolution image
restoration,''
in \textit{Proc. IEEE/CVF Conf. Comput. Vis. Pattern Recognit. (CVPR)},
2022, pp.~5728--5739.

\bibitem{zamir2020mirnet}
S.~W.~Zamir, A.~Arora, S.~Khan, M.~H.~Khan, F.~Shahbaz~Khan,
and L.~Shao,
``Learning enriched features for real image restoration and
enhancement,''
in \textit{Proc. Eur. Conf. Comput. Vis. (ECCV)}, 2020,
pp.~492--511.

\bibitem{hu2018senet}
J.~Hu, L.~Shen, and G.~Sun,
``Squeeze-and-excitation networks,''
in \textit{Proc. IEEE/CVF Conf. Comput. Vis. Pattern Recognit. (CVPR)},
2018, pp.~7132--7141.

\bibitem{sun2021recseg}
L.~Sun, Z.~Fan, Y.~Ding, X.~Huang, and J.~Paisley,
``Joint CS-MRI reconstruction and segmentation with a unified deep
network,''
in \textit{Proc. Int. Conf. Inf. Process. Med. Imag. (IPMI)},
2019, pp.~492--504.

\bibitem{kickingereder2016radiomics}
P.~Kickingereder \textit{et al.},
``Radiomic profiling of glioblastoma: Identifying an imaging
predictor of patient survival with improved performance over
established clinical and radiologic risk models,''
\textit{Radiology}, vol.~280, no.~3, pp.~880--889, Sep.\ 2016.

\bibitem{hou2019pathologygan}
L.~Hou, A.~Agarwal, D.~Samaras, T.~M.~Kurc, R.~R.~Gupta, and
J.~H.~Saltz,
``Robust histopathology image analysis: To label or to synthesize?''
in \textit{Proc. IEEE/CVF Conf. Comput. Vis. Pattern Recognit. (CVPR)},
2019, pp.~8533--8542.

\bibitem{vaswani2017attention}
A.~Vaswani, N.~Shazeer, N.~Parmar, J.~Uszkoreit, L.~Jones,
A.~N.~Gomez, {\L}.~Kaiser, and I.~Polosukhin,
``Attention is all you need,''
in \textit{Adv. Neural Inf. Process. Syst. (NeurIPS)}, vol.~30,
2017, pp.~5998--6008.

\end{thebibliography}

\end{document}
